% Research contribution: all-index Lerch unfolding and a Bessel scaling regime.
% Prepared against manuscript commit 28357e8ca63dd78327db91d9be239d75e4462879.
% This is a section fragment; root manuscript supplies theorem environments.

\section{A Bessel scaling law when the Laurent index follows the derivative order}
\label{bessel:sec:main}

The fixed-index concentration theorem does not cover $n$ comparable with
$k$. In the adjacent-index regime $n=k+r$, with fixed integer $r$, the
zeros close to $a=1$ have a different scale and a different universal
profile. Their distance from one is of order $k^{-2}$, the profile is a
Bessel function, and the Lerch deformation changes these zeros only by an
exponentially small amount. The result holds uniformly up to the classical
endpoint $\rho=1$.

Throughout, let
\[
 F_{n,k}^{\rho}(a)=(-1)^{n+k}\mathcal D_{n,k}^{\rho}(a),\qquad
 f_{n,k}(a)=F_{n,k}^{0}(a).
\]
The exact coefficient and shift identities are
\begin{equation}\label{bessel:eq:shift}
 f_{n,k}(a)=\frac{n!}{a^{k+1}}[t^n](1+t)_k e^{-t\log a},
 \qquad F_{n,k}^{\rho}(a)=\sum_{m\geq0}\rho^m f_{n,k}(a+m).
\end{equation}
The shift series converges absolutely, also at $\rho=1$, for $k\geq1$.
Here $(1+t)_k=\prod_{j=1}^k(t+j)$, so the coefficient polynomial in
$\log a$ has integer coefficients. It is not a reciprocal-gamma Appell
polynomial.

\subsection{An exponentially small spectral tail}

The following estimate retains the angle in Cauchy's coefficient bound.
Replacing $|j+t|$ by $j+|t|$ before optimizing loses the useful decay.

\begin{lemma}[A contour estimate in the adjacent-index range]
\label{bessel:lem:tail}
For integers $n\geq0$, $k\geq1$, $a\geq3/4$ and $0\leq\rho\leq1$,
\begin{equation}\label{bessel:eq:tail}
 \frac{|F_{n,k}^{\rho}(a)-f_{n,k}(a)|}{n!}
 \leq e^7\left(\frac{k}{2}\right)^{k-n}e^{-k/20}
       \left(\frac47+\frac2k\right).
\end{equation}
More generally, for fixed nonnegative integers $p,q$, derivatives
$\partial_\rho^p\partial_a^q$ of the spectral tail satisfy
\[
 \frac{\bigl|\partial_\rho^p\partial_a^q
   (F_{n,k}^{\rho}-f_{n,k})(a)\bigr|}{n!}
 \leq C_p\left(\frac{k+q}{2}\right)^{k+q-n}e^{-(k+q)/20}
\]
for all sufficiently large $k$ (depending on $p$), uniformly for
$a\geq3/4$ and $\rho\in[0,1]$. At $\rho=1$ these are one-sided
derivatives when needed.
\end{lemma}

\begin{proof}
Put $R=k/2$, $t=Re^{i\theta}$ and $v=\cos\theta$. By the arithmetic--
geometric mean inequality,
\begin{align*}
 \frac{|(1+t)_k|}{R^k}
 &\leq\left(\frac{1}{kR^2}
       \sum_{j=1}^k(j^2+R^2+2jRv)\right)^{k/2}\\
 &=(A_k+B_kv)^{k/2},\\
 A_k&=\frac73+\frac2k+\frac{2}{3k^2},\qquad
 B_k=2+\frac2k.
\end{align*}
For $b\geq b_0=7/4$ and $L=\log b$, Cauchy's bound in
\eqref{bessel:eq:shift} therefore gives
\begin{equation}\label{bessel:eq:pointtail}
 \frac{|f_{n,k}(b)|}{n!}
 \leq R^{k-n}b^{-1}
 \exp\left\{k\max_{-1\leq v\leq1}
 \left(\frac12\log(A_k+B_kv)-L(1+v/2)\right)\right\}.
\end{equation}
Write $L_0=\log(7/4)$. Since $1+v/2\geq1/2$, replacing $L$ by
$L_0$ costs at most $-(L-L_0)/2$ in the displayed exponent. Moreover,
\[
 \frac12\log(A_k+B_kv)
 \leq\frac12\log(7/3+2v)+\frac7k.
\]
Indeed $7/3+2v\geq1/3$, while the nonnegative increment is at most
$4/k+2/(3k^2)$; apply $\log(1+x)\leq x$.
If $y=7/3+2v$, maximizing over all $y>0$ gives the upper bound
\[
 \frac12\log\frac2{L_0}-\frac12-\frac5{12}L_0< -\frac1{20}.
\]
For an entirely elementary verification, use $L_0>11/20$ and
$\log(40/11)<13/10$. The expression is then smaller than
$13/20-1/2-11/48=-19/240<-1/20$. These logarithm bounds follow,
for example, by the positive atanh series for the logarithm with a
geometric remainder.
Consequently
\[
 \frac{|f_{n,k}(b)|}{n!}
 \leq e^7R^{k-n}e^{-k/20}b^{-1}
       (b_0/b)^{k/2}.
\]
The function $b^{-1}(b_0/b)^{k/2}$ decreases. Summing it over
$b=a+m$, $m\geq1$, and comparing the tail with its integral gives
at most $b_0^{-1}+2/k=4/7+2/k$, which proves
\eqref{bessel:eq:tail}.

For derivatives, $\partial_a^qf_{n,k}=(-1)^qf_{n,k+q}$, and
$\partial_\rho^p\rho^m=m^{\underline p}\rho^{m-p}$. The latter
factor is bounded by $m^p\leq(a+m)^p$ on the relevant range. For
$K=k+q>2p$, the resulting scalar sum is bounded by
\[
 b_0^{p-1}+\frac{b_0^p}{K/2-p}.
\]
The same estimate proves uniform convergence of the differentiated
series and hence justifies the derivatives, including at $\rho=1$.
\end{proof}

\subsection{The limiting entire profile and its complete expansion}

Fix an integer $r$, and take $k$ large enough that $n=k+r\geq0$. Set
\begin{equation}\label{bessel:eq:profile}
 a=e^{z/k^2},\qquad
 \mathcal H_{r,k}^{\rho}(z)
   =\frac{k^{2r}a^{k+1}}{(k+r)!}F_{k+r,k}^{\rho}(a).
\end{equation}
The entire function defining the limit is
\begin{equation}\label{bessel:eq:Psi}
 \Psi_r(z)=\sum_{j\geq\max(0,-r)}
       \frac{(-z)^{r+j}}{2^j j!(r+j)!}.
\end{equation}
For $z>0$, the standard series for the Bessel function gives
\begin{equation}\label{bessel:eq:Psibessel}
 \Psi_r(z)=(-1)^r(2z)^{r/2}J_r(\sqrt{2z}).
\end{equation}
The series \eqref{bessel:eq:Psi}, rather than the separate powers in
\eqref{bessel:eq:Psibessel}, defines the value at zero and the entire
continuation. Since $J_{-s}=(-1)^sJ_s$ for integer $s$, its positive
zeros are $j_{|r|,m}^2/2$, $m\geq1$. They are simple. The standard
Bessel series and its zero properties are recorded in DLMF
\S\S10.2 and 10.21 \cite{DLMFbessel}; these classical facts are not new claims here.

\begin{theorem}[Uniform Bessel profile]
\label{bessel:thm:profile}
For fixed $r\in\mathbb Z$, on every compact real $z$ interval,
\begin{equation}\label{bessel:eq:profileexp}
 \mathcal H_{r,k}^{\rho}(z)
 =\Psi_r(z)+\frac{1}{3k}\left((2r+5)z\Psi_r'(z)
       +(z-r(2r+5))\Psi_r(z)\right)+O(k^{-2}),
\end{equation}
uniformly for $\rho\in[0,1]$, including every fixed number of $z$
derivatives. There is a complete expansion in integral powers of
$k^{-1}$. Every coefficient in that expansion is independent of $\rho$.
The difference between the profile at any $\rho$ and at zero is
$O(k^r e^{-k/20})$, with the same bound after any fixed number of
$\rho$ derivatives and a fixed number of $z$ derivatives.
\end{theorem}

\begin{proof}
Let $E_{j,k}=e_j(1,\ldots,k)$ be the elementary symmetric functions
of the integers $1,\ldots,k$. Reversing the coefficients of the
rising factorial gives
\[
 (1+t)_k=\sum_{j=0}^kE_{j,k}t^{k-j}.
\]
Consequently the elementary profile is exactly
\begin{equation}\label{bessel:eq:elementaryprofile}
 \mathcal H_{r,k}^{0}(z)
  =\sum_{j=\max(0,-r)}^k\frac{E_{j,k}}{k^{2j}}
       \frac{(-z)^{r+j}}{(r+j)!}.
\end{equation}
Consider the auxiliary polynomial
\[
 E_k(t)=\sum_{j=0}^k\frac{E_{j,k}}{k^{2j}}t^j
       =\prod_{\ell=1}^k\left(1+\frac{\ell t}{k^2}\right).
\]
Its logarithm has the convergent local expansion
\begin{align}\label{bessel:eq:Eexp}
 \log E_k(t)
 &=\sum_{m\geq1}\frac{(-1)^{m-1}t^m}{m k^{2m}}
                \sum_{\ell=1}^k\ell^m\notag\\
 &=\frac t2+\frac{t/2-t^2/6}{k}
       +\frac{-t^2/4+t^3/12}{k^2}+O(k^{-3}).
\end{align}
The remainder is uniform on each fixed complex $t$ disk. For any
desired order, the terms with large $m$ have a uniform geometric
bound because $|\ell t/k^2|\leq |t|/k$. Each retained power sum is
an exact polynomial in $k$, giving a complete expansion. Exponentiating,
\begin{equation}\label{bessel:eq:Efirst}
 E_k(t)=e^{t/2}\left(1+\frac{t/2-t^2/6}{k}
               +\frac{-t^2/8+t^4/72}{k^2}+O(k^{-3})\right).
\end{equation}
Cauchy's coefficient estimate on a fixed circle bounds each
coefficient remainder by $C k^{-N-1}T^{-j}$. Multiplying by
$(-z)^{r+j}/(r+j)!$ and summing proves a locally uniform expansion
of \eqref{bessel:eq:elementaryprofile}, with all derivatives. This
step controls the infinite coefficient sum; coefficientwise convergence
alone would not suffice.

Put $D=z\partial_z-r$. The coefficient of $k^{-1}$ is
\[
 \frac{5D-2D^2}{3}\Psi_r.
\]
The series \eqref{bessel:eq:Psi} satisfies
\begin{equation}\label{bessel:eq:Psiode}
 z\Psi_r''+(1-r)\Psi_r'+\frac12\Psi_r=0.
\end{equation}
Using this equation reduces the last differential expression to the
coefficient displayed in \eqref{bessel:eq:profileexp}. The second
coefficient, useful below, is
\begin{equation}\label{bessel:eq:Hsecond}
 \left(\frac29D(D-1)(D-2)(D-3)-\frac12D(D-1)\right)\Psi_r.
\end{equation}
The remaining coefficients follow by the same finite process from
\eqref{bessel:eq:Eexp}.

Finally, on a compact $z$ interval, $a=e^{z/k^2}\geq3/4$ for
large $k$, and $a^{k+1}$ is uniformly bounded. The spectral-tail
lemma, multiplied by $k^{2r}$, gives the error
$O(k^r e^{-k/20})$. Its differentiated bounds give the same result
for fixed mixed derivatives: each $z$ derivative introduces factors
of $k^{-2}$ when differentiating $a$, while an additional $a$
derivative of the tail costs only a factor $O(k)$. Derivatives of
$a^{k+1}$ cost $O(k^{-1})$ in the $z$ coordinate. Thus all these
errors remain exponentially small and affect none of the algebraic
coefficients.
\end{proof}

\subsection{Simple zeros and deformation independence to every algebraic order}

\begin{theorem}[Bessel locations of adjacent-index zeros]
\label{bessel:thm:zeros}
Fix $r\in\mathbb Z$ and $m\geq1$, and put
$\lambda=j_{|r|,m}^2/2$. For all sufficiently large $k$, every
$\rho\in[0,1]$ has exactly one zero $a_{r,m,k}^{\rho}$ in a
fixed sufficiently small neighborhood of $k^2\log a=\lambda$.
The zero is positive and simple. Uniformly in $\rho$,
\begin{align}\label{bessel:eq:zerosz}
 k^2\log a_{r,m,k}^{\rho}
 &=\lambda-\frac{(2r+5)\lambda}{3k}
       +\frac{\lambda(4r^2+15r+\lambda+20)}{9k^2}
       +O(k^{-3}),\\
 a_{r,m,k}^{\rho}
 &=1+\frac\lambda{k^2}-\frac{(2r+5)\lambda}{3k^3}
       +\frac{\lambda(8r^2+30r+11\lambda+40)}{18k^4}
       +O(k^{-5}).\label{bessel:eq:zerosa}
\end{align}
Both formulas extend to all orders by a recursively computable
expansion with coefficients independent of $\rho$. In addition,
\begin{equation}\label{bessel:eq:deformation}
 \sup_{0\leq\rho\leq1}
 |a_{r,m,k}^{\rho}-a_{r,m,k}^{0}|
       =O(k^{r-2}e^{-k/20}).
\end{equation}
For every fixed positive integer $p$, the same bound holds for
$|\partial_\rho^p a_{r,m,k}^{\rho}|$, once $k$ is sufficiently large.
No assertion about the total positive-zero count follows from this
local scaling theorem.
\end{theorem}

\begin{proof}
The positive Bessel zero $\lambda$ is simple. Choose a fixed small
closed interval about it on which $\Psi_r'$ never vanishes and
whose endpoints have opposite signs. Uniform $C^1$ convergence in
Theorem~\ref{bessel:thm:profile} gives exactly one simple real zero
there for all sufficiently large $k$, uniformly in $\rho$.

Substitute $z=\lambda+d/k+e/k^2$ into the profile expansion. The
$k^{-1}$ coefficient gives
$d=-(2r+5)\lambda/3$. At a zero of $\Psi_r$, equation
\eqref{bessel:eq:Psiode} gives
\[
 \frac{\Psi_r''(\lambda)}{\Psi_r'(\lambda)}
       =\frac{r-1}{\lambda}.
\]
Writing $H_1$ for the $k^{-1}$ coefficient in the profile, direct
differentiation and the same differential equation give
$H_1'(\lambda)/\Psi_r'(\lambda)=\lambda/3$.
Formula~\eqref{bessel:eq:Hsecond} gives
\[
 \frac{H_2(\lambda)}{\Psi_r'(\lambda)}
 =-\frac{2r^3\lambda}{9}-\frac{4r^2\lambda}{3}
   +\frac{2r\lambda^2}{9}-\frac{35r\lambda}{18}
   +\frac{4\lambda^2}{9}-\frac{5\lambda}{6}.
\]
Therefore
\[
 e=-\frac{\Psi_r''(\lambda)d^2/2
           +H_1'(\lambda)d+H_2(\lambda)}{\Psi_r'(\lambda)}
   =\frac{\lambda(4r^2+15r+\lambda+20)}9.
\]
This proves \eqref{bessel:eq:zerosz}; exponentiation gives
\eqref{bessel:eq:zerosa}. At every further order the next unknown
coefficient has the same nonzero multiplier $\Psi_r'(\lambda)$.
The complete profile expansion and a uniform lower bound for the
derivative turn this formal recursion into an asymptotic expansion
with the stated remainders, by the mean value theorem.

The profile difference between $\rho$ and zero is
$O(k^r e^{-k/20})$. Dividing by the uniform derivative lower bound
gives the same error in $z$. The map $a=e^{z/k^2}$ supplies the
additional factor $k^{-2}$ in \eqref{bessel:eq:deformation}.
For derivatives in $\rho$, use the mixed-derivative bounds of the
tail lemma and differentiate the implicit equation. Every positive
$\rho$ derivative of the profile is exponentially small; its
$z$ derivative stays bounded away from zero. Induction proves the
same bound for every fixed root derivative. The required derivatives
at $\rho=1$ are justified by the uniformly convergent differentiated
shift series, for sufficiently large $k$.
\end{proof}

For $r=0$, the limiting function is $J_0(\sqrt{2z})$. The first
zero thus satisfies, for every Lerch parameter including one,
\[
 a=1+\frac{j_{0,1}^2}{2k^2}
       -\frac{5j_{0,1}^2}{6k^3}+O(k^{-4}).
\]
In particular this locates a simple zero of $\gamma_k^{(k)}$ near
one. More generally, the classical specialization of the profile theorem
can be written without the sign-normalized family:
\[
 \frac{k^{2r}e^{(k+1)z/k^2}}{(k+r)!}
 \gamma_{k+r}^{(k)}(e^{z/k^2})
 \longrightarrow (2z)^{r/2}J_r(\sqrt{2z}),\qquad z>0.
\]
For $r>0$ the limiting profile also has a zero of multiplicity $r$
at $z=0$. The theorem deliberately treats the positive Bessel zeros,
not this multiple central zero. Its finer unfolding is sensitive to
the exponentially small spectral tail and is not determined by the
algebraic expansion.

\section{The full small-Lerch-parameter unfolding at arbitrary derivative order}
\label{unfold:sec:main}

The earlier report \cite{LerchBaseline} proves the small-$\rho$ count for every
$n$ at $k=1$. The following generalization handles every pair $(n,k)$.
It also explains why a parity-only count does not extend to arbitrary
derivative order: an explicit rational polynomial evaluated at $\log2$
controls the answer.

\begin{theorem}[All-pair small-$\rho$ count]
\label{unfold:thm:count}
Fix integers $n\geq1$ and $k\geq1$.
If $n\leq k+1$, then $F_{n,k}^{\rho}$ has exactly $n$ simple
positive zeros for all sufficiently small $\rho\geq0$.

If $n>k$, put $r=n-k$ and $B_{n,k}=f_{n,k}(2)$. Then
$B_{n,k}\ne0$, and, for all sufficiently small $\rho>0$, the
number of positive zeros is exactly
\begin{equation}\label{unfold:eq:counts}
 \begin{cases}
 k+1,&r\text{ odd},\\
 k+2,&r\text{ even and }B_{n,k}<0,\\
 k,&r\text{ even and }B_{n,k}>0.
 \end{cases}
\end{equation}
All these zeros are simple. The $k$ noncentral zeros continue
analytically from the elementary endpoint. Every real central branch
has
\begin{equation}\label{unfold:eq:puiseux}
 a(\rho)=1+u\rho^{1/r}+O(\rho^{2/r}),\qquad
 u^r=-\frac{(-1)^r r!}{n!}B_{n,k}.
\end{equation}
The real solutions of the displayed equation give exactly the real
central branches. When $r=1$, the first and second assertions agree.
\end{theorem}

\begin{proof}
At $\rho=0$, the polynomial in $x=-\log a$ is a positive scalar
multiple of
$\prod_{j=1}^k(1+\partial_x/j)x^n$. Each application of
$1+c\partial_x$, $c>0$, preserves real roots, decreases every
inherited multiple-root multiplicity by one, and creates simple
roots strictly between consecutive distinct roots and to the left
of the least root. Hence, for $n>k$, there are $k$ distinct
negative simple roots and a root at zero of multiplicity exactly
$r=n-k$. For $n\leq k+1$ every root is simple. The corresponding
noncentral positive $a$ roots persist analytically for small $\rho$.

The polynomial
\[
 n![t^n](1+t)_k e^{-Lt}\in\mathbb Z[L]
\]
is nonzero (its leading term is $k!(-L)^n$). By Lindemann's
classical theorem \cite{Lindemann}, $\log2$ is transcendental: if it were nonzero
algebraic, its exponential could not equal the algebraic number two.
It cannot be a zero of this nonzero integer polynomial. Thus
$B_{n,k}\ne0$. This arithmetic argument applies to the elementary
coefficient polynomial; it would not apply to a polynomial with
uncontrolled gamma or zeta constants as coefficients.

Writing $a=1+\delta$ and using
$\log^n(1+\delta)/(1+\delta)=\delta^n+O(\delta^{n+1})$ gives
\[
 F_{n,k}^{\rho}(1+\delta)
 =\frac{(-1)^r n!}{r!}\delta^r+\rho B_{n,k}
    +O(\delta^{r+1}+\rho|\delta|+\rho^2).
\]
Set $\rho=\varepsilon^r$ and $\delta=\varepsilon u$.
After division by $\varepsilon^r$, the limiting equation is
\eqref{unfold:eq:puiseux}. Its $r$ complex roots are distinct and
nonzero. The analytic implicit function theorem gives $r$ branches
analytic in $\varepsilon$; Rouch\'e's theorem on a small fixed
$\delta$ disk shows that these exhaust the roots tending to one.
For real $\varepsilon$, branches with real initial values remain
real and the others remain nonreal. For odd $r$ there is one real
root. For even $r$, there are two real roots precisely when the
right-hand side is positive, that is, when $B_{n,k}<0$; otherwise
there are none. Each real branch is simple for small positive
$\varepsilon$.

To exclude further positive roots, take $A$ larger than the largest
elementary root. All $f_{n,k}(a+m)$ have the common sign $(-1)^n$
for $a>A$, so there are no roots there. For sufficiently small $a$,
the positive divergent elementary term dominates the remaining
series, uniformly for $0\leq\rho\leq1/2$. On the compact set
left after removing the elementary root neighborhoods, uniform
convergence to the nonzero elementary function excludes zeros.
This proves both global completeness and the count.
\end{proof}

\begin{corollary}[A necessary condition for uniform positive-parameter saturation]
For $1\leq k\leq n-3$, saturation by $n$ positive zeros cannot hold
for every $\rho\in(0,1]$. At the borderline $k=n-2$, saturation
holds for all sufficiently small positive $\rho$ if and only if
$f_{n,n-2}(2)<0$.
\end{corollary}
\begin{proof}
For $r\geq3$, the count in \eqref{unfold:eq:counts} is strictly
smaller than $k+r=n$. For $r=2$, it is $n$ precisely for the
negative sign, and otherwise equals $n-2$.
\end{proof}

The following are finite exact-arithmetic consequences, separate from
the all-pair analytic theorem. Each sign is certified by a rational
interval for $f_{n,k}(2)/n!$ in the supplied
\texttt{evidence/exact-signs.json}; the code uses the integer rising
factorial and a rational logarithm enclosure.
\begin{center}
\begin{tabular}{rrrr}
\hline
$n$&$k$&Sign of $f_{n,k}(2)$&Count for sufficiently small $\rho>0$\\
\hline
4&2&$+$&2\\
5&3&$+$&3\\
6&2&$+$&2\\
6&4&$-$&6\\
7&5&$-$&7\\
8&6&$+$&6\\
\hline
\end{tabular}
\end{center}
The interval of sufficiently small positive parameters depends on
$(n,k)$ and is not numerically specified by this table. These counts
make no assertion about an intermediate Lerch parameter or the
classical endpoint. The separate full-Lerch Bessel diagnostics in the
package are floating-point observations and are not used in any sign
certificate or global count.

\paragraph{Further problems made precise.}
The central multiple zero in the Bessel scaling theorem lies on a
smaller, exponentially sensitive scale. Its joint $k,\rho$ unfolding,
including possible cancellations between spectral terms, is a natural
next problem. A second problem is a uniform expansion when $|n-k|$
grows, which should connect the fixed-order Bessel regime to a
large-order Bessel or Airy regime. Neither follows from the present
fixed-$r$ theorem. A third problem is the oscillatory sign pattern of
$f_{n,k}(2)$ when $n$ and $k$ grow together; the all-pair small-$\rho$
classification reduces that question to an explicit finite polynomial,
but does not establish an asymptotic sign law.

% Suggested bibliography entries:
% NIST DLMF, Sections 10.2 and 10.21, https://dlmf.nist.gov/10.2,
% https://dlmf.nist.gov/10.21 (series and simple positive Bessel zeros).
% F. Lindemann, Ueber die Zahl pi, Math. Ann. 20 (1882), 213--225,
% DOI 10.1007/BF01446522 (nonzero algebraic logarithms of algebraic numbers).
% Prior incoming report: docs/reports/lerch-zero-bifurcations/article/
% lerch_zero_bifurcations.tex, Section 6 (all-n small-rho theorem at k=1).
